% Ad Familiares 12.2 --- headnote
% ad-familiares-12-02, between 18 September and 5 October 44 BC, Rome

\textit{Cicero to C. Cassius, from Rome between 18 September and
5 October 44 BC --- Perseus dateline \textit{Scr. Romae inter
xiii K. et iii Non. Oct. a. 710 (44)}. The First Philippic has been
delivered (2 September); Antony has answered with the savage senatorial
invective of 19 September; Cicero is now writing to Cassius in his
province of Syria with the temperature of the city in his hand. The
gladiator-metaphor of \textit{Fam.} 12.22 is here in its full form ---
Antony is the \textit{gladiator} hunting for an opening to start the
killing --- and Cicero gives Cassius the names: Piso, who opened the
attacks on 1 August with no one rising to support him; himself, on
2 September; and P. Servilius following. Three consulars who cannot
now come to the Senate in safety.}

\textit{The middle paragraphs are pure Cicero in private vein: a quick
catalogue of disappointments. Cassius's \textit{tuus necessarius} (his
brother-in-law Lepidus) is taken up with his new marriage tie to
Antony and is bursting at the applause for the elder Cassius's
\textit{ludi Apollinares}; the ``other in-law'' (Servilius Isauricus,
or possibly L. Cornelius Lentulus) has been bought off by fresh
``memoranda of Caesar''; and worst of all, somebody in the Senate ---
unnamed, almost certainly Servilius Isauricus --- expects his own son
to be consul in the year that should have been Cassius and Brutus's,
and openly fawns on Antony to that end. The close pivots back to
hope: the consuls-elect, Pansa and Hirtius; Cotta gone to ground in
fatal despair; L. Caesar ill, Sulpicius absent. The counsellors of
the public deliberation are few. ``All hope is in you,'' Cicero
writes, in a sentence that he had already written in May (\textit{Fam.}
12.1) and would write again before the year was out.}
